% ===========================================================================
\section{Full UAIR$^+$ arithmetization}\label{sec:uair}
% ===========================================================================

We assemble the building blocks of \Cref{sec:arith} into a single
UAIR$^+$ instance. We use the \emph{shift-access} extension of the main
paper (Remark~8.2 of the Zinc+ paper), and the
\emph{affine-combination lookup} extension (Remark~8.3).

\subsection{Trace shape and the chained-compression layout}\label{sec:trace-shape}

Fix a number of compressions $N := \mathrm{NUM\_COMPRESSIONS}$ and a
number of rows per compression
$\mathrm{ROWS\_PER\_COMP} := 68$, organized as $4$ \emph{init-prefix
rows} followed by $64$ \emph{round-update output rows}. The trace runs
$N$ chained compressions packed contiguously, with compression
$i \in [0,N)$ occupying rows
$[\,68i,\;68(i+1)\,)$, plus a $4$-row $H_N$ output prefix at rows
$[\,68N,\;68N+4\,)$.
The total number of \emph{active} rows is therefore
$\mathrm{ACTIVE\_ROWS} := 68 N + 4$, and the trace length is
$n := 2^{\nu}$ for the smallest $\nu$ with
$2^{\nu} \ge \mathrm{ACTIVE\_ROWS}$.
The reference implementation uses $N=7$, giving
$\mathrm{ACTIVE\_ROWS} = 480 \le 512 = 2^9$.

Per compression $i$, with $\mathrm{start} := 68i$, four windows are
relevant:
\begin{center}
\begin{tabular}{lll}
\toprule
\textbf{Window} & \textbf{Range} & \textbf{Purpose} \\
\midrule
Init prefix     & $[0, 4)$  & Pinned to $H_i$ via $S_{\mathrm{INIT-PREFIX}}$ \\
Schedule active & $[\mathrm{start},\,\mathrm{start}+48)$ & Anchors of C\ref{con:c7}\\ & & \quad producing $W[k+16]$ for $k+16 \in [16,64)$ \\
Update active   & $[\mathrm{start},\,\mathrm{start}+64)$ & Anchors of C\ref{con:c8}, C\ref{con:c9} \\
Junction window & $[\mathrm{start}+64,\,\mathrm{start}+68)$ & Read by $S_{\mathrm{FF}}$\\ & & \quad for the $H_{i+1} = H_i + (\dots)$ feed-forward \\
\bottomrule
\end{tabular}
\end{center}

\subsection{Column layout}\label{sec:cols}

Following the conventions of UAIR$^+$ in the main paper (see
Definition~8.1 there), each witness vector is partitioned by its
ambient ring. Our index has $\mathbf{q} = ()$ (empty prime tuple),
i.e.\ all witnesses and constraints live in the $\Q[X]$ branch.
Inside the $\Q[X]$ branch, the trace cells are typed as either
\emph{binary-polynomial} (lookup-constrained to $\bitpoly{32}$) or
\emph{integer} (lookup-constrained to a small integer range or to
$\Q$).

Tables \ref{tab:bin-cols} and~\ref{tab:int-cols} list the $20$
binary-polynomial columns and $9$ integer columns of our trace,
respectively, separated into their public and witness portions. We
write $\bp{x}\rowidx{t}$ for the entry of column $\bp{x}$ at row $t$,
and $\bp{x}^{\downarrow\Delta}\rowidx{t} := \bp{x}\rowidx{t+\Delta}$
for shifted access (Remark~8.2 of the main paper), with the convention
$\bp{x}^{\downarrow\Delta}\rowidx{t} = 0$ if $t+\Delta > n$.

\begin{table}[h]
\centering\small
\begin{tabularx}{\linewidth}{>{\raggedright\arraybackslash}p{0.24\linewidth} >{\raggedright\arraybackslash}X >{\raggedright\arraybackslash}p{0.30\linewidth}}
\toprule
\textbf{Column} & \textbf{Role} & \textbf{Lookup set $\subseteq$ Domain} \\
\midrule
\multicolumn{3}{l}{\textit{Public binary-polynomial columns ($9$)}} \\[1pt]
$\col{PA\_A}$           & Boundary column for the $a$-register; carries $H_i$ at init prefixes, $H_N$ at output prefix & $\bitpoly{32} \subseteq \Q^{<32}[X]$ \\
$\col{PA\_E}$           & Boundary column for the $e$-register & $\bitpoly{32} \subseteq \Q^{<32}[X]$ \\
$\col{PA\_OV\_SIG0}$    & Per-coefficient overflow witness for the $\Sigma_0$ rotation lift \eqref{eq:sig0-Q} & $\bitpoly{32} \subseteq \Q^{<32}[X]$ \\
$\col{PA\_OV\_SIG1}$    & Same, for $\Sigma_1$ & $\bitpoly{32} \subseteq \Q^{<32}[X]$ \\
$\col{PA\_OV\_LSIG0}$   & Same, for $\sigma_0$ \eqref{eq:lsig0-Q} & $\bitpoly{32} \subseteq \Q^{<32}[X]$ \\
$\col{PA\_OV\_LSIG1}$   & Same, for $\sigma_1$ & $\bitpoly{32} \subseteq \Q^{<32}[X]$ \\
$\col{PA\_R\_CH2\_COMP}$ & Tail compensator $c_{\mathrm{ch}_2}$ for $r_{\mathrm{ch}_2}$ \eqref{eq:compensators} & $\bitpoly{32} \subseteq \Q^{<32}[X]$ \\
$\col{PA\_R\_MAJ\_COMP}$ & Tail compensator $c_{\mathrm{maj}}$ for $r_{\mathrm{maj}}$ \eqref{eq:compensators} & $\bitpoly{32} \subseteq \Q^{<32}[X]$ \\
$\col{PA\_M}$           & Public message-block words (Table~9 row 77) & $\bitpoly{32} \subseteq \Q^{<32}[X]$ \\
\midrule
\multicolumn{3}{l}{\textit{Witness binary-polynomial columns ($11$)}} \\[1pt]
$\bp{W_a}$              & Working $a$-register $a[t]$ & $\bitpoly{32} \subseteq \Q^{<32}[X]$ \\
$\bp{W_{\Sigma_0}}$     & $\Sigma_0(a[t])$ & $\bitpoly{32} \subseteq \Q^{<32}[X]$ \\
$\bp{W_e}$              & Working $e$-register $e[t]$ & $\bitpoly{32} \subseteq \Q^{<32}[X]$ \\
$\bp{W_{\Sigma_1}}$     & $\Sigma_1(e[t])$ & $\bitpoly{32} \subseteq \Q^{<32}[X]$ \\
$\bp{W_W}$              & Message-schedule word $W[t]$ & $\bitpoly{32} \subseteq \Q^{<32}[X]$ \\
$\bp{W_{\sigma_0}}$     & $\sigma_0(W[t-15])$ & $\bitpoly{32} \subseteq \Q^{<32}[X]$ \\
$\bp{W_{\sigma_1}}$     & $\sigma_1(W[t-2])$ & $\bitpoly{32} \subseteq \Q^{<32}[X]$ \\
$\bp{W_{u_{ef}}}$       & $\Ch$ operand split $u_{ef} = e \AND f$ & $\bitpoly{32} \subseteq \Q^{<32}[X]$ \\
$\bp{W_{u_{\neg e,g}}}$ & $\Ch$ operand split $u_{\neg e,g} = (\NOT e) \AND g$ & $\bitpoly{32} \subseteq \Q^{<32}[X]$ \\
$\bp{W_{\Maj}}$         & $\Maj(a[t],b[t],c[t])$ & $\bitpoly{32} \subseteq \Q^{<32}[X]$ \\
$\bp{W_{\mu\text{-packed}}}$ & Packed integer-carry column (Remark~\ref{rem:mu-budget}) & $\bitpoly{32} \subseteq \Q^{<32}[X]$ \\
\bottomrule
\end{tabularx}
\caption{Binary-polynomial columns ($20$ total: $9$ public, $11$ witness). The $11$ \emph{witness} columns are range-confined to $\bitpoly{32}$ in-circuit by the booleanity sumcheck. The $9$ \emph{public} columns are not committed and are not part of the booleanity sumcheck.}
\label{tab:bin-cols}
\end{table}

\begin{table}[h]
\centering\small
\begin{tabularx}{\linewidth}{>{\raggedright\arraybackslash}p{0.24\linewidth} >{\raggedright\arraybackslash}X >{\raggedright\arraybackslash}p{0.30\linewidth}}
\toprule
\textbf{Column} & \textbf{Role} & \textbf{Lookup set $\subseteq$ Domain} \\
\midrule
\multicolumn{3}{l}{\textit{Public integer columns ($9$ — entire integer section is public)}} \\[1pt]
$S_{\mathrm{INIT-PREFIX}}$  & Selector: $1$ on each of the $N+1$ init-prefix blocks, else $0$ & $\{0,1\} \subseteq \Q$ \\
$S_{\mathrm{FF}}$           & Selector: $1$ on the $N$ junction windows, else $0$ & $\{0,1\} \subseteq \Q$ \\
$S_{\mathrm{MSG-INIT}}$     & Selector for the message-init pin (C\ref{con:msg-init}) & $\{0,1\} \subseteq \Q$ \\
$\col{PA\_K}$               & SHA-256 round constants $K_t$ & $[0..2^{32}-1] \subseteq \Q$ \\
$\col{PA\_C\_C7}$           & Compensator for C\ref{con:c7}; zero on schedule-active rows, arbitrary on inactive rows (verifier-checked, see \Cref{cons:public-checks}) & $\Q$ \\
$\col{PA\_C\_C8}$           & Compensator for C\ref{con:c8}; zero on update-active rows, arbitrary on inactive rows (verifier-checked) & $\Q$ \\
$\col{PA\_C\_C9}$           & Compensator for C\ref{con:c9}; zero on update-active rows, arbitrary on inactive rows (verifier-checked) & $\Q$ \\
$\col{PA\_C\_FF\_A}$        & Compensator for C\ref{con:ff-a}; zero on junction-window rows, arbitrary on inactive rows (verifier-checked) & $\Q$ \\
$\col{PA\_C\_FF\_E}$        & Compensator for C\ref{con:ff-e}; zero on junction-window rows, arbitrary on inactive rows (verifier-checked) & $\Q$ \\
\bottomrule
\end{tabularx}
\caption{Integer columns ($9$ total, all public). Compensator columns ($\col{PA\_C\_*}$) are inspected by the verifier and checked to be zero on the appropriate entries (see below).}
\label{tab:int-cols}
\end{table}


\begin{remark}[What is a compensator column?]\label{rem:what-is-comp}
A \emph{compensator column} is a public column whose sole purpose
is to absorb whatever value the rest of an algebraic constraint
takes on \emph{rows where the constraint does not, in general, hold}. Their goal is to avoid using multiplicative selector polynomials, which would turn linear constraints into non-linear constraints. Compensator columns hinder zero-knowledge, and for this reason in a future iteration of the implementation, we plan to use a Whirlaway-like approach (\url{https://github.com/leanEthereum/leanMultisig/blob/main/minimal_zkVM.pdf}) to forward column shifts, which provides a a (potentially) more zero-knowledge friendly alternative to using multiplicative selector polynomials.


Two distinct families of compensators appear in the SHA UAIR$^+$;
both follow the same pattern but have different types and different
``out-of-scope'' regions.

\paragraph{Active-row compensators
$\col{PA\_C\_C7}, \col{PA\_C\_C8}, \col{PA\_C\_C9}, \col{PA\_C\_FF\_A}, \col{PA\_C\_FF\_E}$
($\Q$-valued).}
Each of constraints \ref{con:c7}--\ref{con:c9} (the SHA recurrence)
and \ref{con:ff-a}--\ref{con:ff-e} (the feed-forward) is of the form
$\mathrm{inner}(t) + \col{pa\_c}(t) \in (X-2)$, where
$\mathrm{inner}(t)$ is the SHA-256 algebraic expression that
\emph{must} vanish modulo $X-2$ on the constraint's active row range
(schedule-active anchors, update-active anchors, junction-window rows)
and is \emph{arbitrary} elsewhere. The compensator $\col{pa\_c}(t)$
takes value $0$ on the active rows and absorbs $-\mathrm{inner}(t)$
on the inactive rows so that the overall expression lies in $(X-2)$
\emph{everywhere}, sidestepping a multiplicative selector that would
inflate the constraint degree. Why $\Q$-valued: on inactive rows,
$\mathrm{inner}(2)$ can be any rational, so the compensator must
range over all of $\Q$. The active-row zero-pinning is enforced by
the verifier directly on \texttt{public\_trace} (\Cref{cons:public-checks}).

\paragraph{Tail compensators
$\col{PA\_R\_CH2\_COMP}, \col{PA\_R\_MAJ\_COMP}$ (bit-polynomial-typed).}
The virtual residuals
$r_{\mathrm{ch}_2}, r_{\mathrm{maj}}$ of \Cref{sec:ch-maj} are
booleanity-pinned to $\{0,1\}^{32}$ per coefficient. On every
\emph{interior} row they collapse to clean
$\XOR$/$\Maj$-$\XOR$ patterns and lie in $\bitpoly{32}$
automatically; on the last two rows the length-$2$ forward shifts
read off-trace zero-padding and the bit pattern \emph{slips outside}
$\{0,1\}^{32}$ unless corrected. The tail compensators
$c_{\mathrm{ch}_2}, c_{\mathrm{maj}}$ carry the exact boundary
bit-pattern that restores the residual to $\bitpoly{32}$ on those
two tail rows (and are zero everywhere else; see \eqref{eq:compensators}).


In both families, the compensator is \emph{public}: the verifier
either supplies the values directly (active-row case, with the
zero-on-active rows discharged by direct row-wise inspection) or
reads them from \texttt{public\_trace} (tail case, with the
$\bitpoly{32}$ structural property type-enforced). No in-circuit
constraint pins compensator values explicitly; their structure is
discharged outside the algebraic part of the protocol.
\end{remark}

\begin{remark}[Public boundaries: all init-prefix blocks live on $\col{PA\_A}/\col{PA\_E}$]
\label{rem:boundaries}
The columns $\col{PA\_A}$ and $\col{PA\_E}$ play a dual role.
On each init-prefix block
$[\,68i,\,68i+4\,)$ they carry $H_i = (h_0^{(i)},\dots,h_7^{(i)})$
in the SHA shift-register convention (so the four-row block encodes
$(d,c,b,a)$ in $\col{PA\_A}$ and $(h,g,f,e)$ in $\col{PA\_E}$). On the
junction window $[\,68i+64,\,68i+68\,)$ they carry $H_i$ again, this
time as the addend in the feed-forward sum. The verifier supplies both
blocks, so $H_i$ is duplicated on disjoint row windows to keep both the
init-prefix pin and the feed-forward compensator at degree~$1$ in the
trace MLEs. The terminal block $[\,68N,\,68N+4\,)$ carries $H_N$.
\end{remark}

\subsection{Constraints}\label{sec:constraints}

We organize the constraints by family. All algebraic constraints below
are $\Q[X]$-constraints; the relevant ideal is one of $(X^{32}-1)$ or
$(X-2)$, or the zero ideal $(0)$. We follow the Zinc+ convention of
labelling algebraic constraints C1, C2, \dots\ for ease of
cross-reference with the implementation.

The constraint family is supplemented by a small set of
\emph{verifier-side public-column structural checks}
(\Cref{cons:public-checks} below) that the verifier discharges by
direct inspection of \texttt{public\_trace} before the algebraic
checks begin. These structural checks pin the public compensators
and compensators row-wise without consuming any selector column or
any algebraic constraint.

\phantomsection\label{cons:public-checks}
\paragraph{Public-column properties.}

We assume that the public input is well formed, in the following sense (i.e.\ we do not enforce the properties below in our constraints). 
\begin{itemize}[leftmargin=2em]
  \item Every public binary-polynomial column lies in $\bitpoly{32}$
    coefficient-wise. Concretely: for each row $t$ of every
    public-section column listed in \Cref{tab:bin-cols}, every
    coefficient of the cell is in $\{0,1\}$.
  \item Every cell of every public integer column lies in its
    declared lookup set of \Cref{tab:int-cols}, namely
    $\{0,1\}$ for the $S$-selectors, $[0..2^{32}-1]$ for $\col{PA\_K}$,
    and any $\Q$-value for the compensators.
  \item Checked by the verifier: Compensator columns are zero on
    their constraint's row set $\mathcal{T}$ (defined in the row-set
    table just below):
    \begin{align*}
      \col{PA\_C\_C7}\rowidx{t}    &= 0 \quad\text{for every schedule-active row } t, \\
      \col{PA\_C\_C8}\rowidx{t}    &= 0 \quad\text{for every update-active row } t, \\
      \col{PA\_C\_C9}\rowidx{t}    &= 0 \quad\text{for every update-active row } t, \\
      \col{PA\_C\_FF\_A}\rowidx{t} &= 0 \quad\text{for every junction-window row } t, \\
      \col{PA\_C\_FF\_E}\rowidx{t} &= 0 \quad\text{for every junction-window row } t.
    \end{align*}
  \item Checked by the verifier: Tail-compensator zero-on-inner-rows and boundary closed form:
    for every $t$ with $t+2 < n$,
    \[
      \col{PA\_R\_CH2\_COMP}\rowidx{t} = 0,
      \qquad
      \col{PA\_R\_MAJ\_COMP}\rowidx{t} = 0,
    \]
    and on $t \in \{n-2, n-1\}$ the closed-form values \eqref{eq:compensators}.
  \item Round-constant pinning: $\col{PA\_K}\rowidx{t}$ equals the
    SHA-256 round constant $K_{t'}$ for $t' = t + 3$ on every
    update-active row $t$ (i.e.\ on the C\ref{con:c8}, C\ref{con:c9}
    re-anchored row variable).
  \item Selector value pattern: each of $S_{\mathrm{INIT-PREFIX}},
    S_{\mathrm{FF}}, S_{\mathrm{MSG-INIT}}$ takes the value $1$ on
    its declared support set (\Cref{tab:int-cols}) and $0$ elsewhere.
\end{itemize}
These are $O(n)$ row-wise inspections of \texttt{public\_trace} that
the verifier already iterates through to derive public lifted
evaluations; the marginal cost is negligible.

\paragraph{Row-variable convention.}\label{sec:row-conv}

Every algebraic constraint below is a row-local expression in the
trace cells, instantiated at every row in a specified subset
$\mathcal{T} \subseteq [\,0, n\,)$. We use a single row variable
$t$ throughout. Within an expression, every cell access is either
unshifted ($\bp{x}\rowidx{t}$) or a forward shift
$\bp{x}^{\downarrow\Delta}\rowidx{t} = \bp{x}\rowidx{t+\Delta}$
for $\Delta > 0$, with the off-trace convention
$\bp{x}\rowidx{t+\Delta} = 0$ when $t + \Delta \ge n$. Backward
shifts never appear: where the natural SHA-spec form reads
``$W[t'-16]$'' or ``$a[t'-3]$'', we re-index by setting $t := t'-16$
(or $t := t'-3$) so all reads land at $t + \Delta$ with
$\Delta \ge 0$. The re-indexed $t$ then ranges over a window
shifted by the corresponding offset relative to the SHA round
index $t'$.

\smallskip

The row sets $\mathcal{T}$ are reused across constraint families:
\begin{center}
\small
\begin{tabular}{lll}
\toprule
\textbf{Name} & \textbf{Row set $\mathcal{T}$} & \textbf{Used by} \\
\midrule
Active rows           & $[0, \mathrm{ACTIVE\_ROWS})$ & C\ref{con:c1}, C\ref{con:c2} \\
Schedule-active rows  & $\bigcup_{i=0}^{N-1} [\,68i,\,68i + 48\,)$ & C\ref{con:c4}, C\ref{con:c6}, C\ref{con:c7} \\
Update-active rows    & $\bigcup_{i=0}^{N-1} [\,68i,\,68i + 64\,)$ & C\ref{con:c8}, C\ref{con:c9} \\
Init-prefix rows      & $\bigcup_{i=0}^{N}   [\,68i,\,68i + 4\,)$  & C\ref{con:c10}, C\ref{con:c11} \\
Junction-window rows  & $\bigcup_{i=0}^{N-1} [\,68i + 64,\,68(i{+}1)\,)$ & C\ref{con:ff-a}, C\ref{con:ff-e} \\
Message-seed rows     & $\bigcup_{i=0}^{N-1} [\,68i,\,68i + 16\,)$ & C\ref{con:msg-init} \\
Every row             & $[0, n)$                                  & C\ref{con:mu-zero}, C\ref{con:lkup-uef}, C\ref{con:lkup-uneg}, C\ref{con:lkup-maj} \\
\bottomrule
\end{tabular}
\end{center}

A $(X-2)$-ideal constraint listed for ``schedule-active rows''
(say) holds on those rows by virtue of the SHA recurrence; on
\emph{inactive} rows a public $\Q$-valued compensator
$\col{PA\_C\_*}$ absorbs whatever value the inner expression takes
so that the $(X-2)$-ideal membership holds on every row of the
trace, not just on $\mathcal{T}$ — see \Cref{rem:what-is-comp}.
The compensator is required to be zero on $\mathcal{T}$ itself
(verifier-checked, \Cref{cons:public-checks}), which is what makes
the constraint genuinely enforce the SHA recurrence on $\mathcal{T}$.

\paragraph{Rotation constraints.}

Following \Cref{lem:big-sigma}, for every \emph{active row} $t$:
\begin{align}
\bp{W_a}\rowidx{t} \cdot \rho_{0}(X) - \bp{W_{\Sigma_0}}\rowidx{t} - 2\,\col{PA\_OV\_SIG0}\rowidx{t} &\in (X^{32}-1), \tag{C1}\label{con:c1}\\
\bp{W_e}\rowidx{t} \cdot \rho_{1}(X) - \bp{W_{\Sigma_1}}\rowidx{t} - 2\,\col{PA\_OV\_SIG1}\rowidx{t} &\in (X^{32}-1). \tag{C2}\label{con:c2}
\end{align}

Following \Cref{lem:small-sigma}, with $\SHR^{r}(\bp{W_W})$ realized
as a virtual column over $\bp{W_W}$, for every \emph{schedule-active row}
$t$ (i.e.\ $t \in \bigcup_i [\,68i,\,68i+48\,)$; equivalently, the
re-indexed input row corresponding to SHA round $t' = t + 16$):
\begin{align}
\bp{W_W}^{\downarrow 1}\rowidx{t} \cdot \rho_{\sigma_0}(X) + \SHR^{3}(\bp{W_W}^{\downarrow 1})\rowidx{t}
   - \bp{W_{\sigma_0}}^{\downarrow 1}\rowidx{t} - 2\,\col{PA\_OV\_LSIG0}^{\downarrow 1}\rowidx{t} &\in (X^{32}-1), \tag{C4}\label{con:c4}\\
\bp{W_W}^{\downarrow 14}\rowidx{t} \cdot \rho_{\sigma_1}(X) + \SHR^{10}(\bp{W_W}^{\downarrow 14})\rowidx{t}
   - \bp{W_{\sigma_1}}^{\downarrow 14}\rowidx{t} - 2\,\col{PA\_OV\_LSIG1}^{\downarrow 14}\rowidx{t} &\in (X^{32}-1). \tag{C6}\label{con:c6}
\end{align}
The forward-shift offsets ($1$ and $14$) align the SHA-spec
$\sigma_0(W[t'-15])$ and $\sigma_1(W[t'-2])$ reads (with
$t' = t + 16$ the SHA round index) onto cells of the
$\bp{W_{\sigma_0}}, \bp{W_{\sigma_1}}$ columns at $t + 1$, $t + 14$
respectively.

\paragraph{Message-schedule modular sum.}

The SHA spec's message-schedule recurrence
$W[t'] = W[t'-16] + \sigma_0(W[t'-15]) + W[t'-7] + \sigma_1(W[t'-2])$
(SHA round index $t'$) is re-indexed by $t := t' - 16$ so all reads
become forward shifts of $t$. The constraint is then enforced on
every \emph{schedule-active row} $t$:
\begin{equation}\tag{C7}\label{con:c7}
\begin{aligned}
\bp{W_W}^{\downarrow 16}\rowidx{t}
  - \bp{W_W}\rowidx{t}
  - \bp{W_{\sigma_0}}^{\downarrow 1}\rowidx{t}
  - \bp{W_W}^{\downarrow 9}\rowidx{t}
  &- \bp{W_{\sigma_1}}^{\downarrow 14}\rowidx{t} \\
  + 2X^{31} \cdot \mu_W[t] + \col{PA\_C\_C7}\rowidx{t} &\in (X-2),
\end{aligned}
\end{equation}
where $\mu_W[t] := \mathrm{ShiftR}(\bp{W_{\mu\text{-packed}}}, 0, 2)\rowidx{t}$
extracts bits $0$--$1$ of the packed carry column. Off
schedule-active rows the compensator $\col{PA\_C\_C7}$ absorbs the
inner expression so the $(X-2)$-membership holds vacuously
(\Cref{rem:what-is-comp}).

\paragraph{Register update for $a$.}

The SHA round-update for register $a$ at SHA round index $t'$ reads
from registers $a[t'-1], \ldots, a[t'-4]$ and $e[t'-1], \ldots, e[t'-4]$
(via the shift-register convention $h[t'] = e[t'-3]$, etc.; see
\Cref{rem:shift-register}). Re-indexing by $t := t' - 3$ shifts the
deepest backward reads onto the row variable $t$ itself; the output
row $t' = t + 3$ becomes a $\downarrow 4$ access. The constraint is
enforced on every \emph{update-active row} $t$:
\begin{equation}\tag{C8}\label{con:c8}
\begin{aligned}
&\bp{W_a}^{\downarrow 4}\rowidx{t}
  - \bp{W_e}\rowidx{t}
  - \bp{W_{\Sigma_1}}^{\downarrow 3}\rowidx{t}
  - \bp{W_{u_{ef}}}^{\downarrow 3}\rowidx{t}
  - \bp{W_{u_{\neg e,g}}}^{\downarrow 3}\rowidx{t} \\
&\quad - \col{PA\_K}^{\downarrow 3}\rowidx{t}
  - \bp{W_W}^{\downarrow 3}\rowidx{t}
  - \bp{W_{\Sigma_0}}^{\downarrow 3}\rowidx{t}
  - \bp{W_{\Maj}}^{\downarrow 3}\rowidx{t} \\
&\qquad + 2X^{31} \cdot \mu_a[t] + \col{PA\_C\_C8}\rowidx{t} \in (X-2),
\end{aligned}
\end{equation}
where $\mu_a[t] := \mathrm{ShiftR}(\bp{W_{\mu\text{-packed}}}, 2, 3)\rowidx{t}$.
$\col{PA\_C\_C8}$ absorbs the inner expression off update-active
rows.

\paragraph{Register update for $e$.}

Same re-indexing as C\ref{con:c8}: the constraint is enforced on
every update-active row $t$ (with $t' = t + 3$ the SHA round
index):
\begin{equation}\tag{C9}\label{con:c9}
\begin{aligned}
&\bp{W_e}^{\downarrow 4}\rowidx{t}
  - \bp{W_a}\rowidx{t}
  - \bp{W_e}\rowidx{t}
  - \bp{W_{\Sigma_1}}^{\downarrow 3}\rowidx{t}
  - \bp{W_{u_{ef}}}^{\downarrow 3}\rowidx{t}
  - \bp{W_{u_{\neg e,g}}}^{\downarrow 3}\rowidx{t} \\
&\quad - \col{PA\_K}^{\downarrow 3}\rowidx{t}
  - \bp{W_W}^{\downarrow 3}\rowidx{t}
  + 2X^{31} \cdot \mu_e[t] + \col{PA\_C\_C9}\rowidx{t} \in (X-2).
\end{aligned}
\end{equation}

\paragraph{Init-prefix pinning.}

For every init-prefix row $t$ across all $N+1$ blocks (including the
$H_N$ output prefix):
\begin{align}
S_{\mathrm{INIT-PREFIX}}\rowidx{t} \cdot \bigl(\bp{W_a}\rowidx{t} - \col{PA\_A}\rowidx{t}\bigr) &= 0, \tag{C10}\label{con:c10}\\
S_{\mathrm{INIT-PREFIX}}\rowidx{t} \cdot \bigl(\bp{W_e}\rowidx{t} - \col{PA\_E}\rowidx{t}\bigr) &= 0. \tag{C11}\label{con:c11}
\end{align}
Both are constraints in the zero ideal (strict equality), gated by a
boolean selector. Because $S_{\mathrm{INIT-PREFIX}}$ is a public
$\{0,1\}$-valued column, the multiplication preserves degree~$1$ in
the trace MLEs only on a per-row basis; the constraint is degree-$2$
overall, but it carries the zero ideal so it does not enter the
projected sum-check's effective degree count.

\paragraph{Feed-forward to the next compression.}

For every \emph{junction-window row} $t$ (i.e.\
$t \in \bigcup_i [\,68i + 64,\,68(i + 1)\,)$):
\begin{align}
\bp{W_a}^{\downarrow 4}\rowidx{t} - \bp{W_a}\rowidx{t} - \col{PA\_A}\rowidx{t}
   + 2X^{31} \cdot \mu_{ff,a}[t] + \col{PA\_C\_FF\_A}\rowidx{t}
   &\in (X-2), \tag{C12}\label{con:ff-a}\\
\bp{W_e}^{\downarrow 4}\rowidx{t} - \bp{W_e}\rowidx{t} - \col{PA\_E}\rowidx{t}
   + 2X^{31} \cdot \mu_{ff,e}[t] + \col{PA\_C\_FF\_E}\rowidx{t}
   &\in (X-2), \tag{C13}\label{con:ff-e}
\end{align}
with $\mu_{ff,a}[t] := \mathrm{ShiftR}(\bp{W_{\mu\text{-packed}}}, 8, 1)\rowidx{t}$
and $\mu_{ff,e}[t] := \mathrm{ShiftR}(\bp{W_{\mu\text{-packed}}}, 9, 1)\rowidx{t}$.
The compensators $\col{PA\_C\_FF\_A}, \col{PA\_C\_FF\_E}$ absorb
the inner expressions off the junction window.

The Table~9 (from the Zinc+ main paper) materialization constraints C13--C15 of the spec are
\emph{absent} from this UAIR; their role is taken by the virtual
binary-polynomial residuals \eqref{eq:rch1}--\eqref{eq:rmaj},
pinned by the booleanity sumcheck via the closing-override mechanism.
We retain the C-numbering jump for consistency with the implementation.

\paragraph{Message-init pinning.}

For every row $t$ in any of the $N$ message-seed windows
$[\,68i,\,68i+16\,)$:
\begin{equation}\tag{C16}\label{con:msg-init}
S_{\mathrm{MSG-INIT}}\rowidx{t} \cdot \bigl(\bp{W_W}\rowidx{t} - \col{PA\_M}\rowidx{t}\bigr) = 0.
\end{equation}
This pins the first $16$ message-schedule cells of each compression to
the public message word column.

\paragraph{Range-pin for the unused bits of the packed carry column.}
\begin{equation}\tag{C17}\label{con:mu-zero}
\mathrm{ShiftR}(\bp{W_{\mu\text{-packed}}}, 10, 22)\rowidx{t} = 0
\quad \text{for every row } t.
\end{equation}
Equivalently, the high $22$ bits of $\bp{W_{\mu\text{-packed}}}$ are
zero; the low $10$ bits are read by the five $\mathrm{ShiftR}$ virtuals
of \Cref{rem:mu-budget}.

\paragraph{Affine-combination lookups for $\Ch$ / $\Maj$.}

In addition to the per-column lookups of \Cref{tab:bin-cols}, we
impose the following $\bitpoly{32}$-membership constraints on
specified affine combinations of trace cells (cf.\ Lemmas~8.9
and~8.11 of the Zinc+ paper, in the \emph{linear-combination lookup}
variant). These are exactly the residuals
$r_{\mathrm{ch}_1}, r_{\mathrm{ch}_2}, r_{\mathrm{maj}}$ defined in
\Cref{sec:ch-maj}, restated here in the $\downarrow$-notation of the
constraint section. With $\bp{x}^{\downarrow\Delta}\rowidx{t} = \bp{x}\rowidx{t+\Delta}$
(zero past the trace tail) and the public bit-poly tail compensators
of \eqref{eq:compensators}:
\begin{align}
\bp{W_e}^{\downarrow 2}\rowidx{t} + \bp{W_e}^{\downarrow 1}\rowidx{t} - 2\,\bp{W_{u_{ef}}}^{\downarrow 2}\rowidx{t} &\in \bitpoly{32}, \tag{C21}\label{con:lkup-uef}\\
\bp{W_e}^{\downarrow 2}\rowidx{t} - \bp{W_e}\rowidx{t} + 2\,\bp{W_{u_{\neg e,g}}}^{\downarrow 2}\rowidx{t} + 2\,\col{PA\_R\_CH2\_COMP}\rowidx{t} &\in \bitpoly{32}, \tag{C22}\label{con:lkup-uneg}\\
\bp{W_a}\rowidx{t} + \bp{W_a}^{\downarrow 1}\rowidx{t} + \bp{W_a}^{\downarrow 2}\rowidx{t} - 2\,\bp{W_{\Maj}}^{\downarrow 2}\rowidx{t} - 2\,\col{PA\_R\_MAJ\_COMP}\rowidx{t} &\in \bitpoly{32}, \tag{C23}\label{con:lkup-maj}
\end{align}
on every row $t$. The implementation registers the three left-hand
sides as virtual binary-polynomial columns (no commit), and the
booleanity sumcheck enforces \ref{con:lkup-uef}--\ref{con:lkup-maj}
coefficient-wise on every row simultaneously.

\ref{con:lkup-uef} needs no compensator: at the tail rows
$t \in \{n-2, n-1\}$ the off-trace shift values reduce the residual
to $\bp{W_e}^{\downarrow 1}\rowidx{n-2} = \bp{W_e}\rowidx{n-1}$ (resp.\ $0$),
already in $\{0,1\}$ per coefficient. \ref{con:lkup-uneg} and
\ref{con:lkup-maj} \emph{do} need their compensators
$\col{PA\_R\_CH2\_COMP}, \col{PA\_R\_MAJ\_COMP}$ on those two tail
rows; without them the residuals slip to $\{-1, 0\}$
(\ref{con:lkup-uneg}) or $\{0, 1, 2\}$ (\ref{con:lkup-maj}) per
coefficient — see the discussion preceding \eqref{eq:compensators}.

